% -------------------------------------------------------------------
\section{Results and Discussion}
\label{sec:results}
% -------------------------------------------------------------------

We evaluate the proposed system over 120-second Mininet experiments under three conditions: (1)~\emph{Baseline}---standard TCP with \texttt{pfifo} queueing and no ECN; (2)~\emph{Binary ECN}---binary ECN toggle using an XGBoost model at threshold 0.83, following the approach of~\cite{welzl2024}; (3)~\emph{Proportional ECN}---our proportional RED threshold adjustment using LightGBM at threshold 0.56. Network conditions: 10\,Mbit/s bottleneck, 30\,ms RTT, 6 background flows. Both TCP Reno and TCP Cubic are evaluated as the foreground flow.

\subsection{Model Classification Performance}

Table~\ref{tab:models} summarizes classification performance across all trained models. LightGBM consistently outperforms all other classifiers on the combined 57,808-row dataset, achieving F1\,=\,0.9723, ROC-AUC\,=\,0.9965, and PR-AUC\,=\,0.9960. This strong performance is attributable to three design decisions: (i) the leak-free multi-signal labeling strategy producing a well-separated class boundary; (ii) 71 candidate features (45 retained after selection) providing rich temporal and relational context; and (iii) 7-fold TimeSeriesSplit with per-fold threshold calibration ensuring robust temporal generalization.

\subsection{TCP Reno --- Mininet Experiment Results}

Table~\ref{tab:reno} presents the Reno results. Our proportional LightGBM approach achieves 1.60\,Mbps throughput---a 15.5\% improvement over the 1.39\,Mbps baseline and 7.4\% over the binary ECN approach (1.49\,Mbps). Average CWND increases from 5.3 to 6.0 segments, indicating that proportional ECN signaling sustains a larger congestion window by avoiding abrupt halving caused by packet drops.

The proportional mechanism proactively injects 18 ECN signals, sustaining a higher average CWND (6.0 vs.\ 5.3 segments) and operating closer to bottleneck capacity. At this higher utilization point, the increased CWND oscillation ($\sigma$: 2.37$\rightarrow$2.56) is an expected side effect---a deliberate trade-off that yields a net $+15.5\%$ throughput gain. The binary ECN approach at threshold 0.83 fires zero ECN signals under these moderate congestion conditions, demonstrating the importance of proper threshold calibration.

\begin{table}[!t]
\renewcommand{\arraystretch}{1.3}
\caption{TCP Reno Results (120\,s, 10\,Mbit/s, 30\,ms RTT)}
\label{tab:reno}
\centering
\begin{tabular}{lccc}
\toprule
\textbf{Metric} & \textbf{Baseline} & \textbf{Binary ECN} & \textbf{Prop. ECN} \\
\midrule
Throughput (Mbps)       & 1.39  & 1.49  & 1.60 (+15.5\%) \\
Avg CWND (segs)         & 5.27  & 5.71  & 6.00 \\
CWND Stability ($\sigma$) & 2.37 & 2.38 & 2.56 \\
Avg RTT (ms)            & 44.9  & 45.7  & 44.7 \\
ECN Signals Sent        & ---   & 0     & 18 \\
\bottomrule
\multicolumn{4}{p{0.95\columnwidth}}{%
$\dagger$ Retransmission counts are from the final \texttt{ss -ti} snapshot.
The iperf3 sender report recorded \textbf{1,077 cumulative retransmits} for
the Reno baseline over 120\,s, confirming significant packet loss under
baseline conditions. Both ECN approaches recorded 0 iperf3 retransmits.}
\end{tabular}
\end{table}

% Fig. 2 — Reno CWND trace
\begin{figure}[!t]
\centering
\includegraphics[width=\columnwidth]{cwnd_trace(Reno).png}
\caption{TCP Reno CWND traces under three conditions (120\,s, 10\,Mbit/s, 30\,ms RTT). The proportional ECN approach maintains a higher sustained CWND ($\sim$6.0 segments) through 18 proactive ECN signals, while baseline shows frequent deep drops.}
\label{fig:cwnd_reno}
\end{figure}

% Fig. 3 — Reno throughput comparison
\begin{figure}[!t]
\centering
\includegraphics[width=\columnwidth]{through(Reno).png}
\caption{Throughput comparison for TCP Reno across Baseline, Binary ECN, and Proportional ECN conditions (120\,s, 10\,Mbit/s, 30\,ms RTT). The proportional approach achieves a 15.5\% improvement over baseline.}
\label{fig:throughput_reno}
\end{figure}

\subsection{TCP Cubic --- Mininet Experiment Results}

Table~\ref{tab:cubic} presents the Cubic results. The most significant finding is the complete elimination of retransmissions: the baseline records 755 retransmissions---characteristic of Cubic's aggressive convex growth phase overflowing the bottleneck queue---while both ECN conditions reduce this to exactly zero.

\begin{table}[!t]
\renewcommand{\arraystretch}{1.3}
\caption{TCP Cubic Results (120\,s, 10\,Mbit/s, 30\,ms RTT)}
\label{tab:cubic}
\centering
\begin{tabular}{lccc}
\toprule
\textbf{Metric} & \textbf{Baseline} & \textbf{Binary ECN} & \textbf{Prop. ECN} \\
\midrule
Throughput (Mbps)         & 1.39  & 1.25  & 1.26 \\
Retransmissions           & 755   & 0 (100\%$\downarrow$) & 0 (100\%$\downarrow$) \\
Avg CWND (segs)           & 5.31  & 4.96  & 5.10 \\
CWND Stability ($\sigma$) & 2.03  & 1.87  & 1.80 (+11.3\%) \\
Avg RTT (ms)              & 45.9  & 47.8  & 48.5 \\
ECN Signals Sent          & ---   & 0     & 2 \\
\bottomrule
\end{tabular}
\end{table}

Our proportional approach achieves the best CWND stability ($\sigma=1.80$ vs.\ 1.87 binary ECN vs.\ 2.03 baseline), an 11.3\% improvement over baseline. Smoother CWND behavior directly translates to reduced jitter and more predictable application latency. The throughput reduction from 1.39\,Mbps (baseline) to 1.26\,Mbps is a well-documented ECN trade-off: preventing queue overflow necessarily keeps CWND below the saturation point, accepting marginally lower throughput in exchange for zero packet loss and improved stability~\cite{aqm_lstm}.

In the Cubic case, even the passive RED+ECN queue marking at the bottleneck is sufficient to prevent overflow. Our system additionally sends 2 targeted signals at highest-confidence prediction points, achieving the best CWND stability across all three conditions.

% Fig. 4 — Cubic CWND trace
\begin{figure}[!t]
\centering
\includegraphics[width=\columnwidth]{cwnd_trace(cubic).png}
\caption{TCP Cubic CWND traces under three conditions. Baseline shows large oscillations (12--13 segments with deep drops) causing 755 retransmissions. The proportional ECN approach achieves the smoothest trace ($\sigma=1.80$, 11.3\% stability improvement) with 2 targeted ECN signals.}
\label{fig:cwnd_cubic}
\end{figure}

\subsection{ECN Signal Activation Analysis}

Table~\ref{tab:ecn_signals} contrasts ECN signal counts and outcomes across all conditions. LightGBM at threshold 0.56 fires 18 signals for Reno and 2 for Cubic; the binary XGBoost at threshold 0.83 fires zero in both cases. This difference reflects threshold calibration: a threshold of 0.83 calibrated on a heavily imbalanced dataset with $\sim$1.2\% positive class rate suppresses activation under moderate congestion, while our balanced dataset (48\% positive) yields a lower, more responsive optimal threshold.

\begin{table}[ht]
\renewcommand{\arraystretch}{1.3}
\caption{ECN Signal Count and Threshold Comparison}
\label{tab:ecn_signals}
\centering
\footnotesize
\begin{tabularx}{\columnwidth}{@{}l l l l L@{}}
\toprule
\textbf{Cond.} & \textbf{Model} & \textbf{Thr.} & \textbf{Sigs.} & \textbf{Key Outcome} \\
\midrule
Reno-Bin. & XGBoost  & 0.83 & 0  & No proactive response; 0 retrans \\
Reno-Prop. & LightGBM & 0.56 & 18 & $+15.5\%$ throughput \\
Cubic-Bin. & XGBoost  & 0.83 & 0  & 0 retrans via passive RED+ECN \\
Cubic-Prop. & LightGBM & 0.56 & 2  & 0 retrans + best CWND ($\sigma=1.80$) \\
\bottomrule
\end{tabularx}
\end{table}

\subsection{Comparative Summary}

% Fig. 5 — Overall system performance
\begin{figure}[!t]
\centering
\includegraphics[width=\columnwidth]{overall.png}
\caption{Overall system performance comparison across TCP Reno and TCP Cubic for Baseline, Binary ECN, and Proportional ECN conditions. Reno demonstrates 15.5\% throughput gain; Cubic demonstrates 100\% retransmission elimination with 11.3\% CWND stability improvement.}
\label{fig:overall}
\end{figure}

\begin{table}[ht]
\renewcommand{\arraystretch}{1.3}
\caption{Overall System Performance}
\label{tab:summary}
\centering
\footnotesize
\begin{tabularx}{\columnwidth}{@{}L L@{}}
\toprule
\textbf{Dimension} & \textbf{Our System} \\
\midrule
Dataset size       & 57,808 rows (combined Reno+Cubic) \\
Models trained     & Single unified Reno+Cubic model \\
Label strategy     & Multi-signal consensus (leak-free) \\
Class balance      & 48.07\% positive / 51.93\% negative \\
Features           & 71 engineered / 45 retained \\
Best model F1      & 0.9723 (LightGBM) \\
ECN mechanism      & \texttt{tc qdisc} proportional RED adjustment \\
Deployment threshold & 0.56 (calibrated) \\
Reno throughput gain & $+15.5\%$ vs.\ baseline \\
Cubic retransmissions & 0 (100\% elimination) + best CWND stability \\
ECN signals (Reno) & 18 proactive signals \\
\bottomrule
\end{tabularx}
\end{table}

Three core findings emerge from the results. First, the labeling strategy is the single most impactful design decision: the multi-signal consensus approach produces a well-separated, naturally balanced class boundary that enables high-fidelity classification (F1\,=\,0.9723) without relying on aggressive oversampling. Second, proportional ECN at the queue-management level delivers superior performance compared to binary packet-marking: for Reno it activates proactively (18 signals) and delivers higher throughput; for Cubic it achieves additional CWND smoothing beyond passive AQM. Third, training a single unified model across both CC variants generalizes successfully---confirmed by the near-identical $\sim$28,900 Reno and Cubic sample counts contributing equally to the F1\,=\,0.9723 result, eliminating the need for per-algorithm model specialization.

% -------------------------------------------------------------------
\section{Conclusion}
\label{sec:conclusion}
% -------------------------------------------------------------------

This paper presented a proactive TCP congestion control system that combines real-time ML-based loss prediction with proportional RED queue management to prevent packet loss before it occurs. Three key contributions distinguish the work: a multi-signal consensus labeling strategy that eliminates data leakage and produces a naturally balanced training set; a single LightGBM model trained on combined Reno and Cubic data with 71 engineered features (45 retained after automated selection); and a proportional artificial ECN mechanism that adjusts RED thresholds continuously as a function of prediction confidence.

Experimental results on Mininet demonstrate the effectiveness of the approach. For TCP Cubic, the system eliminates 100\% of retransmissions ($755 \rightarrow 0$) and achieves an 11.3\% improvement in CWND stability over the uncontrolled baseline. For TCP Reno, a 15.5\% throughput improvement is achieved by sustaining a higher average congestion window through proactive congestion signaling. In both cases, the proportional LightGBM mechanism outperforms the binary ECN approach, demonstrating the value of confidence-proportional queue management.

The work also establishes that a single unified model trained across both Reno and Cubic variants generalizes successfully, removing the need for per-algorithm model specialization. This is a step toward a more broadly applicable congestion prediction system that could extend to additional CC algorithms without retraining.

Several limitations remain. All experiments are conducted in a controlled Mininet environment with a fixed 10\,Mbit/s bottleneck and 30\,ms RTT; real-world networks exhibit greater variability in background traffic patterns, RTT dynamics, and path asymmetry. Future work should evaluate the system across higher-bandwidth paths, cellular and WiFi link types, and variable competing flow counts. Training on actual Internet traces---such as MAWI or CAIDA datasets---would strengthen generalizability claims. Additionally, extending the proportional control mechanism to adjust per-flow pacing rate in addition to RED thresholds may yield further improvements in throughput-stability trade-off management.

% -------------------------------------------------------------------
\begin{thebibliography}{17}

\bibitem{rfc5681}
E.~Blanton, V.~Paxson, and M.~Allman, \textit{TCP Congestion Control}, RFC 5681, Sep. 2009.

\bibitem{rfc9438}
L.~Xu, S.~Ha, I.~Rhee, V.~Goel, and L.~Eggert, \textit{CUBIC for Fast and Long-Distance Networks}, RFC 9438, Aug. 2023.

\bibitem{rfc3168}
S.~Floyd, K.~Ramakrishnan, and D.~Black, \textit{The Addition of Explicit Congestion Notification (ECN) to IP}, RFC 3168, Sep. 2001.

\bibitem{manet}
O.~J.~Adaramola and J.~R.~Olasina, ``Machine Learning-Driven Congestion Prediction in Mobile Ad-Hoc Networks Through Modelling Approaches,'' 2023.

\bibitem{ptcl}
A.~Akhtar, I.~Zafar, M.~A.~Zaki, and M.~Khurram, ``Machine Learning-based fixed access Network Interface Congestion Prediction,'' \textit{Pakistan Journal of Engineering, Technology \& Science}, vol.~12, no.~1, pp.~26--36, 2024.

\bibitem{stgcn}
B.~Magar, R.~Kulkarni, V.~Khatavkar, and H.~Vanjari, ``Predictive network congestion management for enterprise systems: a graph neural network approach with operational insights,'' 2025.

\bibitem{ali2024}
I.~Ali, S.~Hong, and T.~Cheung, ``Congestion or no congestion: Packet loss identification and prediction using machine learning,'' in \textit{Proc. 16th Int. Conf. on Platform Technology and Service (PlatCon)}, 2024, pp.~1--6.

\bibitem{labonne}
M.~Labonne, C.~Chatzinakis, and A.~Olivereau, ``Predicting Bandwidth Utilization on Network Links Using Machine Learning,'' in \textit{Proc. 2020 European Conference on Networks and Communications (EuCNC)}, 2020, pp.~242--247.

\bibitem{wide2025}
Y.~Wang, ``Advanced Network Traffic Prediction Using Deep Learning Techniques: A Comparative Study of SVR, LSTM, GRU, and Bidirectional LSTM Models,'' \textit{ITM Web of Conferences}, vol.~70, 2025.

\bibitem{rl_mesh}
S.~Mahajan, P.~Srideviponmalar, R.~Harikrishnan, and K.~Kotecha, ``Reinforcement Learning Based Congestion Control Technique for Wireless Mesh Networks,'' \textit{International Journal of Networked and Distributed Computing}, 2025.

\bibitem{dqn_tcp}
E.~A\v{g}lamazlar, E.~Eken, and H.~B.~Ge\c{c}ici, ``A Deep Reinforcement Learning-Based TCP Congestion Control Algorithm: Design, Simulation, and Evaluation,'' \textit{arXiv preprint}, 2026.

\bibitem{aqm_lstm}
C.~A.~Gomez, X.~Wang, and A.~Shami, ``Intelligent Active Queue Management Using Explicit Congestion Notification,'' in \textit{Proc. 2019 IEEE Global Communications Conference (GLOBECOM)}, 2019.

\bibitem{remora}
A.~Kumar and N.~Hemrajani, ``Optimized Extreme Gradient Boosting with Remora Algorithm for Congestion Prediction in Transport Layer,'' \textit{International Journal of Computer Network and Information Security}, vol.~16, no.~3, 2024.

\bibitem{ccn}
P.~Bazmi and M.~Keshtgary, ``A Neural Network Based Congestion Control Algorithm for Content-Centric Networks,'' \textit{Journal of Advanced Computer Science \& Technology}, 2014.

\bibitem{wsn}
G.-W.~Lee, S.-Y.~Lee, and E.-N.~Huh, ``Congestion Prediction Modeling for Quality of Service Improvement in Wireless Sensor Networks,'' \textit{Sensors}, vol.~14, no.~5, 2014.

\bibitem{tafa_survey}
Z.~Tafa and V.~Milutinovic, ``Machine Learning in Congestion Control: A Survey on Selected Algorithms and a New Roadmap to their Implementation,'' \textit{arXiv preprint arXiv:2112.15522}, 2021.

\bibitem{welzl2024}
M.~Welzl, S.~Islam, and M.~von Stephanides, ``Real-Time TCP Packet Loss Prediction Using Machine Learning,'' \textit{IEEE Access}, vol.~12, pp.~159622--159634, 2024. DOI: 10.1109/ACCESS.2024.3488511.

\end{thebibliography}

\end{document}
